% Source: 01 Properties_of_Matter 26(1).docx; formatting in layout.tex.
\OutlinedChapterTitle{\ScriptTitle}


\textbf{The students\ldots{}}


\Objective{know the definition of the term \textit{chemistry.}}

\Objective{know the GHS pictograms and their meanings.}

\Objective{understand why experiments are essential in science.}

\Objective{know the basic building blocks of a lab report and know how to write a lab report for an experiment.}

\Objective{\textbf{ know the definition of the particle model and are able to apply it to different situations.}}

\Objective{\textbf{ know what a chemical reaction is, can classify reactions as exothermic or endothermic and can draw the corresponding energy diagrams.}}

\Objective{can formulate reaction equations using names of substances and understand the process on a particle level.}

\Objective{know the definition of the term \textit{substance}}

\Objective{can convert from decimal to scientific to prefix notation and vice versa}

\Objective{can convert units}

\Objective{understand the material properties \textit{melting point }and \textit{boiling point} and know how different factors influence them.}

\Objective{know the characteristics of the states of matter and the transitions between them on the macroscopic level (= our daily lives) and on a particle level.}

\Objective{\textbf{understand the property }\textbf{\textit{solubility}}\textbf{ and know the different terms connected with solubility (solute, solvent, dissolved, saturated solution, solubility).}}

\Objective{understand solubility diagrams and can extract information of a diagram.}

\Objective{understand the material property \textit{density} and know how different factors influence it.}

\Objective{are able to solve simple problems concerning density.}

\Objective{understand the concept which explains the temperature of a substance.}

\Objective{know the term \textit{diffusion} and \textit{osmosis }and can explain them by using examples.}

\Objective{\textbf{ know and understand the terms }\textbf{\textit{homogeneous mixtures}}\textbf{ and }\textbf{\textit{heterogeneous mixtures}}\textbf{.}}

\Objective{know the names of the different types of mixtures and can draw them on a particle level and can name an example.}

\Objective{can name and describe a number of important separation methods.}

\Objective{understand the principles which lie behind these separation methods also on a particle level.}

\Objective{\textbf{ can use the terms }\textbf{\textit{pure substance}}\textbf{, }\textbf{\textit{mixture}}\textbf{, }\textbf{\textit{element}}\textbf{, and }\textbf{\textit{compound }}\textbf{to explain the division of matter which is commonly used in chemistry. }}

\Objective{understand all exercises and experiments}

\ObjectivesNote{Objectives in bold are also relevant for future exams}
